\chapter{Your Python Toolkit and Environment Setup}
\label{ch:python}

\chapterguide%
  {You should be comfortable creating folders and running a command in a terminal.}
  {Create a clean project workspace, manually download the course datasets into dedicated dataset folders, create and activate a virtual environment, install the project dependencies from the included requirements file, and save your practical work as Jupyter notebooks in \texttt{all\_experiments/}.}

Machine learning combines ideas with computation. From this chapter onward, major methods are paired with complete Python examples and small public Kaggle datasets. You will still type, run, modify, and explain the code yourself. The objective of the setup is not automation; it is to make every path predictable enough that the machine-learning ideas stay in focus.

\section{The course project layout}

Use one root folder for the course. Inside it, keep \emph{datasets} and \emph{experiments} separate:

\begin{lstlisting}[style=mlterminal]
UTP_MACHINE_LEARNING/
  all_datasets/                  # one subfolder per dataset
    titanic_dataset/
    medical_cost_personal_dataset/
    breast_cancer_dataset/
    ames_housing_dataset/
    mall_customers_dataset/
    optical_recognition_dataset/
    credit_card_fraud_dataset/
    student_performance_dataset/
    heart_failure_dataset/
    red_wine_quality_dataset/
    bike_sharing_dataset/

  all_experiments/              # 10 canonical practical notebooks
    lab01-foundations_environment_data.ipynb
    lab02-learning_loop_geometry.ipynb
    ...
    lab10-reinforcement_integration.ipynb

  requirements.txt
\end{lstlisting}

\begin{keyidea}
\texttt{all\_datasets/} is only the dataset \emph{root}. Do not throw every CSV into that root. Each downloaded dataset gets its own folder because a Kaggle download may contain one file, several related files, a README, or supporting metadata. Your code should therefore use paths such as \texttt{all\_datasets/titanic\_dataset/Titanic-Dataset.csv}.
\end{keyidea}

\section{Download the course datasets manually from Kaggle}

Open the dataset page by clicking on the name of the dataset in the table below, download the archive, extract it, and place the extracted files in the dedicated folder.

\subsection{Core datasets reused across several chapters}

\begin{mltable}
\footnotesize
\begin{tabularx}{\linewidth}{@{}p{0.38\linewidth}X@{}}
\toprule
\rowcolor{MLPaleBlue!92}
Kaggle dataset & Main use in these notes \\
\midrule
\href{https://www.kaggle.com/datasets/yasserh/titanic-dataset}{Titanic Dataset} & auditing, baselines, model selection, logistic regression, trees \\
\href{https://www.kaggle.com/datasets/mirichoi0218/insurance}{Medical Cost Personal Data} & first pipeline, linear and ridge regression \\
\href{https://www.kaggle.com/datasets/uciml/breast-cancer-wisconsin-data}{Breast Cancer Wisconsin} & geometry, evaluation, KNN, naive Bayes, SVM, limited labels \\
\href{https://www.kaggle.com/datasets/shashanknecrothapa/ames-housing-dataset}{Ames Housing} & random forests and gradient boosting \\
\href{https://www.kaggle.com/datasets/vjchoudhary7/customer-segmentation-tutorial-in-python}{Mall Customers} & unsupervised workflow, K-means, DBSCAN \\
\href{https://www.kaggle.com/datasets/yadavhim/optical-recognition-of-handwritten-digits}{Optical Recognition of Handwritten Digits} & PCA and representation learning \\
\href{https://www.kaggle.com/datasets/miadul/credit-card-fraud-detection-dataset}{Credit Card Fraud Detection} & anomaly detection and imbalanced evaluation \\
\href{https://www.kaggle.com/datasets/larsen0966/student-performance-data-set}{Student Performance} & unfamiliar-data audit, delimiter/schema comparison, stage checkpoint \\
\href{https://www.kaggle.com/datasets/fedesoriano/heart-failure-prediction}{Heart Failure Prediction} & classification evaluation on a new mixed-type schema \\
\href{https://www.kaggle.com/datasets/uciml/red-wine-quality-cortez-et-al-2009}{Red Wine Quality} & ensemble transfer exercise \\
\href{https://www.kaggle.com/datasets/lakshmi25npathi/bike-sharing-dataset}{Bike Sharing} & chronological splitting, leakage checks, regression transfer lab \\
\bottomrule
\end{tabularx}
\end{mltable}

\subsection{One download, step by step}

Use the same procedure for every dataset:

\begin{mlsteps}
  \item Open the Kaggle Dataset link in the table.
  \item Click \textbf{Download}. If Kaggle gives you a ZIP file, extract it normally using your operating system.
  \item Create the dedicated course folder named in the second column if it does not already exist.
  \item Copy the required file or files into that dedicated folder. Keep additional documentation from the download there too if it helps you understand the data.
  \item Confirm that the path in your notebook exactly matches the folder name and filename, including capitalization.
\end{mlsteps}

\begin{pitfall}
A path error is not a machine-learning error. If \texttt{pd.read\_csv(...)} raises \texttt{FileNotFoundError}, first print or inspect your working directory, then compare the actual folder and filename with the path in the notebook. Do not ``fix'' the problem by moving random files into the project root.
\end{pitfall}

\section{Install Python and make it available in the terminal}

Download a current supported release from the \href{https://www.python.org/downloads/}{official Python downloads page}.

\begin{description}[leftmargin=1.15in,style=nextline]
  \item[Windows] Run the official installer, install for your account, and allow it to add Python to \texttt{PATH}. Reopen the terminal, then use \texttt{py}.
  \item[macOS] Run the signed \texttt{.pkg} installer with the default options, then restart Terminal. Do not remove the system Python in \texttt{/usr/bin}.
  \item[Linux] Install Python 3 through your distribution's package manager. It normally adds \texttt{python3} to \texttt{PATH}; do not replace the operating system's Python.
\end{description}

If Windows does not recognize \texttt{py}, add Python's command directory to your system \texttt{PATH} so PowerShell and Command Prompt can find it. Open the Start menu, search for \emph{Edit environment variables for your account}, select \texttt{Path}, click \emph{Edit}, and add the directory shown by the Python installer. With the current install manager, this directory is normally:

\begin{lstlisting}[style=mlterminal]
%LocalAppData%\Python\bin
\end{lstlisting}

On macOS, the installer should configure the path automatically. If \texttt{which -a python3} does not show \texttt{/usr/local/bin/python3}, make sure \texttt{/usr/local/bin} is in the path used by the default Z shell:

\begin{lstlisting}[style=mlterminal]
echo 'export PATH="/usr/local/bin:$PATH"' >> ~/.zshrc
source ~/.zshrc
\end{lstlisting}

For Debian or Ubuntu:

\begin{lstlisting}[style=mlterminal]
sudo apt update
sudo apt install python3 python3-venv python3-pip
\end{lstlisting}

Verify the installation and see which executable the terminal found:

\begin{lstlisting}[style=mlterminal]
# Windows
py --version
py -m pip --version
where py

# macOS or Linux
python3 --version
python3 -m pip --version
which -a python3
\end{lstlisting}

\section{Create an isolated virtual environment}

Python's standard-library \texttt{venv} module creates a lightweight environment with its own installed packages. This keeps the course dependencies separate from other Python projects. \href{https://docs.python.org/3/library/venv.html}{Python \texttt{venv} documentation}.

The repository includes a normal text file named \texttt{requirements.txt} in the project root. It contains the course dependencies:

\begin{lstlisting}[style=mlterminal]
numpy
pandas
matplotlib
scikit-learn
jupyter
\end{lstlisting}

If you installed standard Python, use the command for your operating system to create a virtual environment. The name \texttt{.venv} is a common convention, but you may choose a different name.

\begin{lstlisting}[style=mlterminal]
# Windows
py -m venv .venv

# macOS or Linux
python3 -m venv .venv
\end{lstlisting}

Then staying in the same folder activate the virtual environment PowerShell or Terminal:

\begin{lstlisting}[style=mlterminal]
# Windows
.venv\Scripts\Activate.ps1

# macOS/Linux
source .venv/bin/activate
\end{lstlisting}

After activation, upgrade pip and install the course requirements:

\begin{lstlisting}[style=mlterminal]
# Windows
py -m pip install --upgrade pip
py -m pip install -r requirements.txt

# macOS or Linux
python3 -m pip install --upgrade pip
python3 -m pip install -r requirements.txt
\end{lstlisting}

Package installation is normally a one-time setup. On later study sessions, activate \texttt{.venv} again; you do not need to reinstall the requirements each time.

\begin{keyidea}
Do not install random packages globally whenever a tutorial asks for one. Activate that project environment first, and install into it.
\end{keyidea}

\section{Verify the scientific stack}

Use the practical lab below as the single environment check. It imports the core scientific packages and records their versions so the same notebook also serves as a reproducibility note. If all four imports succeed, the core environment is ready.


\begin{quickcheck}
\begin{enumerate}
  \item Why should datasets, notebooks, and generated code listings live in separate folders?
  \item Why is a project-local virtual environment safer than installing course packages globally?
  \item If a notebook raises \texttt{FileNotFoundError}, what should you inspect before changing any data files?
\end{enumerate}
\end{quickcheck}

\begin{chapterrecap}
\begin{itemize}
  \item \texttt{all\_datasets/} is the dataset root; every Kaggle dataset lives inside its own dedicated subfolder.
  \item A dataset folder may legitimately contain several related files, as with Optical Digits, Student Performance, and Bike Sharing.
  \item Install the included \texttt{requirements.txt} inside a project-local virtual environment.
  \item Save your practical work as notebooks in \texttt{all\_experiments/}; a good notebook records code, plots, results, and interpretation.
\end{itemize}
\end{chapterrecap}
